% Image credits for the photographs used in this volume (Book 1).
% Machine-readable license records live in images/book1/CREDITS.md; keep
% the two files in sync when adding or removing a photograph.
% Arabic edition: personal names are transliterated as in the captions, with
% the original spelling in parentheses (the legally required credit); user
% names, institutions, works and licences stay verbatim.
\chapter*{حقوق الصور}

رسوم هذا الكتاب أصلية. أما الصور الفوتوغرافية فمستعملة بمقتضى
رخصها الحرة، مع الشكر لأصحابها:

\begin{itemize}
  \item \emph{جون دالتون} (فصل الذرات والجزيئات): نقش حفره
        تشارلز تيرنر (Charles Turner) عن جوزيف ألن (Joseph Allen)، 1834،
        Library of Congress عبر Wikimedia Commons، ملك عام.
  \item \emph{جدول دالتون للعناصر، 1808} (فصل الذرات
        والجزيئات): لوحة من كتاب جون دالتون \emph{A New System of Chemical
        Philosophy}، مسحها Internet Archive، عبر Wikimedia Commons،
        ملك عام.
  \item \emph{البوكسيت} (فصل المواد الأولية): صورة التقطها فيرنر
        شيلمان (Werner Schellmann)، عبر Wikimedia Commons، CC BY-SA 2.5.
  \item \emph{الغرانيت} (فصل المواد النقية): صورة التقطها يوريكو
        زيمبريس (Eurico Zimbres)، عبر Wikimedia Commons، CC BY-SA 2.0 br.
  \item \emph{كبريتات النحاس اللامائية} (فصل التعرف على المواد):
        صورة التقطها ف.~أولن (W.~Oelen)، عبر Wikimedia Commons، CC BY-SA 3.0.
  \item \emph{بلورات كبريتات النحاس} (فصل التعرف على المواد):
        صورة التقطها Stephanb، عبر Wikimedia Commons، CC BY-SA 3.0.
  \item \emph{أنطوان لوران لافوازييه وماري آن بولز لافوازييه}
        (فصل المعادلات الموزونة): لوحة رسمها جاك لوي دافيد (Jacques-Louis David)،
        1788، Metropolitan Museum of Art عبر Wikimedia Commons، ملك
        عام.
  \item \emph{هاتف من الباكليت} (فصل البلاستيك): صورة التقطها هولغر
        إلغارد (Holger Ellgaard)، عبر Wikimedia Commons، CC BY-SA 3.0.
  \item \emph{إرنست رذرفورد} (فصل داخل الذرة): صورة فوتوغرافية،
        George Grantham Bain Collection، Library of Congress، عبر
        Wikimedia Commons، ملك عام.
  \item \emph{الهاليت} (فصل الأيونات): صورة التقطها روبرت م.~لافينسكي
        (Robert M.~Lavinsky)، عبر Wikimedia Commons، CC BY-SA 3.0.
  \item \emph{تمثال الحرية} (فصل الفلزات والتآكل):
        صورة التقطها Elcobbola، عبر Wikimedia Commons، ملك عام.
  \item \emph{دميتري مندلييف} و\emph{جدوله لسنة 1869} (فصل الجدول
        الدوري): عبر Wikimedia Commons، ملك عام.
  \item \emph{سديم السرطان} (فصل العناصر): NASA، ESA، ج.~هستر
        (J.~Hester) وأ.~لول (A.~Loll)، ملك عام.
  \item \emph{الصوديوم} (فصل الطبقات الإلكترونية): صورة التقطها Dnn87، عبر
        Wikimedia Commons، CC BY-SA 3.0.
  \item \emph{الكلور في أمبولة} (فصل الطبقات الإلكترونية):
        صورة التقطها ف.~أولن (W.~Oelen)، عبر Wikimedia Commons، CC BY-SA 3.0.
  \item \emph{أميديو أفوغادرو} (فصل المول): طبعة حجرية عن
        رسم للفنان ك.~سانتييه (C.~Sentier)، 1856، Edgar Fahs Smith collection، عبر
        Wikimedia Commons، ملك عام.
  \item \emph{بلورات الثلج} (فصل القطبية): صور التقطها ويلسون
        بنتلي (Wilson Bentley)، 1902، عبر Wikimedia Commons، ملك عام.
  \item \emph{لويس باستور} (فصل الكيمياء الفراغية): صورة التقطها
        بول نادار (Paul Nadar)، عبر Wikimedia Commons، ملك عام.
  \item \emph{بلورات طرطرات الأمونيوم} (فصل الكيمياء الفراغية):
        صورة التقطتها ماري-لان تاي بامار (Marie-Lan Taÿ Pamart) عبر Wikimedia Commons،
        CC BY 4.0.
  \item \emph{سفانته أرهينيوس} (فصل الأحماض والقواعد): صورة محفورة ضوئيًا،
        1909، عبر Wikimedia Commons، ملك عام.
  \item \emph{عمود فولتا} (فصل الخلايا): صورة التقطها GuidoB، عبر
        Wikimedia Commons، CC BY-SA 3.0.
  \item \emph{والاس كاروثرز في مختبره} (فصل البوليمرات):
        مصوّر مجهول، عبر Wikimedia Commons، ملك عام.
\end{itemize}

روابط المصادر الكاملة وعناوين رخصة كل صورة مذكورة في
\texttt{images/book1/CREDITS.md} في مستودع الكتاب.
والرسوم التوضيحية للمخاطر هي رسوم النظام المنسق عالميًا لتصنيف
المواد الكيميائية ووسمها، كما رُسمت لحزمة
\texttt{ghsystem} في \TeX\ Live.

\bigskip
وإلى جانب الصور الفوتوغرافية والرسوم الأصلية، بعض الأشكال
رسوم توضيحية وُلِّدت بالذكاء الاصطناعي، أُنتجت بمولّد الصور لدى OpenAI انطلاقًا من
مطالبات كتبها محررو الكتاب، وروجعت للتأكد من دقتها الكيميائية
قبل إدراجها. والمطالبة الكاملة لكل صورة مولَّدة
مسجَّلة في \texttt{images/book1/ai/PROMPTS.md} في المستودع.
\clearpage
